\documentclass[10pt, letterpaper]{article}

% Packages:
\usepackage[
    ignoreheadfoot,
    top=0.6 cm,
    bottom=0.6 cm,
    left=0.8 cm,
    right=0.8 cm,
]{geometry}
\usepackage{titlesec}
\usepackage{tabularx}
\usepackage{array}
\usepackage[dvipsnames]{xcolor}
\definecolor{primaryColor}{RGB}{0, 0, 0} % Stile monocromatico professionale (B&W)
\usepackage{enumitem}
\usepackage{fontawesome5}
\usepackage{amsmath}
\usepackage{graphicx}
\usepackage[
    pdftitle={Curriculum Vitae - Melania Roccuzzo},
    pdfauthor={Melania Roccuzzo},
    colorlinks=true,
    urlcolor=primaryColor
]{hyperref}
\usepackage{calc}
\usepackage{bookmark}
\usepackage{lastpage}
\usepackage{changepage}
\usepackage{paracol}
\usepackage{ifthen}
\usepackage{needspace}
\usepackage{iftex}

\ifPDFTeX
    \input{glyphtounicode}
    \pdfgentounicode=1
    \usepackage[utf8]{inputenc}
    \usepackage{lmodern}
\fi

\usepackage[italian]{babel} % Sillabazione e regole tipografiche italiane
\usepackage{microtype} % Migliora la giustificazione ed evita testo fuori dai margini

% Settings:
\AtBeginEnvironment{adjustwidth}{\partopsep0pt}
\pagestyle{empty}
\setcounter{secnumdepth}{0}
\setlength{\parindent}{0pt}
\setlength{\topskip}{0pt}
\setlength{\columnsep}{0cm}
\emergencystretch 2em % Impedisce alle righe di oltrepassare il margine destro

\titleformat{\section}{\needspace{4\baselineskip}\bfseries\large}{}{0pt}{}[\vspace{1pt}\titlerule]

\titlespacing{\section}{-1pt}{0.25 cm}{0.12 cm}

\renewcommand\labelitemi{$\circ$}
\newenvironment{highlights}{
    \begin{itemize}[
        topsep=0.06 cm,
        parsep=0.04 cm,
        partopsep=0pt,
        itemsep=0pt,
        leftmargin=0.4 cm + 10pt
    ]
}{
    \end{itemize}
}

\newenvironment{onecolentry}{
    \begin{adjustwidth}{0.00001 cm}{0.00001 cm}
}{
    \end{adjustwidth}
}

\newenvironment{twocolentry}[2][]{
    \onecolentry
    \def\secondColumn{#2}
    \setcolumnwidth{\fill, 6.0 cm}
    \begin{paracol}{2}
}{
    \switchcolumn \raggedleft \secondColumn
    \end{paracol}
    \endonecolentry
}

\newenvironment{header}{
    \setlength{\topsep}{0pt}\par\kern\topsep\centering\linespread{1.1}
}{
    \par\kern\topsep
}

\let\hrefWithoutArrow\href
\renewcommand{\href}[2]{\hrefWithoutArrow{#1}{\ifthenelse{\equal{#2}{}}{ }{#2 }\raisebox{.15ex}{\footnotesize \faExternalLink*}}}

\begin{document}

    \newcommand{\AND}{\unskip
        \cleaders\copy\ANDbox\hskip\wd\ANDbox
        \ignorespaces
    }
    \newsavebox\ANDbox
    \sbox\ANDbox{}

    % HEADER (CON FOTO PROFESSIONALE A DESTRA)
    \noindent
    \begin{minipage}[c]{0.78\textwidth}
        {\fontsize{22 pt}{24 pt}\selectfont\textbf{Melania Roccuzzo}}\par
        \vspace{0.12 cm}
        {\normalsize\textit{Dottoressa in Scienze e Tecnologie Alimentari $\cdot$ Laureanda Magistrale in Nutrizione Umana}}\par
        \vspace{0.2 cm}
        \small
        \mbox{{\color{black}\footnotesize\faMapMarker*}\hspace*{0.13cm}Catania, Italia}%
        \kern 0.45 cm%
        \mbox{\hrefWithoutArrow{mailto:melania.roccuzzo01@email.com}{{\color{black}\footnotesize\faEnvelope[regular]}\hspace*{0.13cm}melania.roccuzzo01@email.com}}%
        \par\vspace{0.1 cm}
        \mbox{\hrefWithoutArrow{tel:+39-3664355965}{{\color{black}\footnotesize\faPhone*}\hspace*{0.13cm}+39 366 4355965}}%
        \kern 0.45 cm%
        \mbox{\hrefWithoutArrow{https://linkedin.com/in/melania-roccuzzo}{{\color{black}\footnotesize\faLinkedinIn}\hspace*{0.13cm}linkedin.com/in/melania-roccuzzo}}%
    \end{minipage}%
    \hfill
    \begin{minipage}[c]{0.20\textwidth}
        \raggedleft
        \IfFileExists{foto.jpg}{%
            \includegraphics[width=2.4cm, height=3.0cm, keepaspectratio]{foto.jpg}%
        }{%
            \IfFileExists{foto.png}{%
                \includegraphics[width=2.4cm, height=3.0cm, keepaspectratio]{foto.png}%
            }{%
                \framebox[2.4cm]{\rule{0pt}{2.8cm}\centering\footnotesize\color{gray}{\begin{tabular}{c}\faUser[regular]\\[0.15cm]Carica\\\texttt{foto.jpg}\end{tabular}}}%
            }%
        }
    \end{minipage}

    \vspace{0.15 cm}

    % PROFILO PROFESSIONALE
    \section{Profilo Professionale}
    \begin{onecolentry}
        Laureata in Scienze e Tecnologie Alimentari e laureanda magistrale in Scienze della Nutrizione Umana, con solida preparazione scientifica nella biochimica degli alimenti, composizione nutrizionale e impatto della nutrizione sul benessere psico-fisico. Forte predisposizione all'ascolto attivo e alla relazione d'aiuto, maturata nella ricerca sui DCA e in contesti operativi a contatto con il pubblico. Ricerco un'opportunità di tirocinio formativo / affiancamento clinico presso studi nutrizionali o poliambulatori a Catania, per supportare con rigore, precisione e proattività la gestione quotidiana dei pazienti, l'elaborazione dei piani nutrizionali e le attività dello studio.
    \end{onecolentry}

    \vspace{0.1 cm}

    % FORMAZIONE ACCADEMICA
    \section{Formazione Accademica}

        \begin{twocolentry}{
            \textit{Catania, Italia}\vspace{0 cm}
            
            \mbox{\textit{2024-In corso (Laurea prevista: 2027)}}
        }
            \textbf{Laurea Magistrale in Scienze della Nutrizione Umana}\vspace{0 cm}
            
            \textit{Università degli Studi di Catania}
        \end{twocolentry}
        \vspace{0.1 cm}
        \begin{onecolentry}
            \begin{highlights}
                \item Percorso avanzato su: biochimica della nutrizione, fisiopatologia, dietoterapia applicata, nutrizione sportiva, valutazione dello stato nutrizionale e fabbisogni nelle varie fasi della vita.
                \item In preparazione per l'Esame di Stato per l'Abilitazione all'esercizio della professione di Biologo Nutrizionista (ONB).
            \end{highlights}
        \end{onecolentry}

        \vspace{0.25 cm}

        \begin{twocolentry}{
            \textit{Catania, Italia}\vspace{0 cm}
            
            \textit{2021-2024}
        }
            \textbf{Laurea Triennale in Scienze e Tecnologie Alimentari}\vspace{0 cm}
            
            \textit{Università degli Studi di Catania}
        \end{twocolentry}
        \vspace{0.1 cm}
        \begin{onecolentry}
            \begin{highlights}
                \item \textbf{Tesi di Laurea:} \textit{``Disturbi Alimentari e Salute Mentale''}.
                \item Condotta un'analisi critica della letteratura scientifica internazionale correlando specifici pattern dietetici al benessere neurochimico e psicologico; sviluppato un quadro metodologico per integrare evidenze scientifiche e proposte pratiche mirate alla prevenzione e al supporto dei disturbi alimentari (DCA).
                \item Approfondite competenze su: composizione molecolare e merceologica degli alimenti, igiene e sicurezza alimentare (HACCP), etichettatura nutrizionale e biodisponibilità dei micro/macronutrienti.
            \end{highlights}
        \end{onecolentry}

    \vspace{0.1 cm}

    % ESPERIENZA & RELAZIONE CON L'UTENTE
    \section{Esperienza Professionale \& Patient Care Skills}

        \begin{twocolentry}{
            \textit{Remoto / Italia}\vspace{0 cm}
            
            \textit{In corso}
        }
            \textbf{Customer Support Specialist}\vspace{0 cm}
            
            \textit{Qomodo (Fintech Startup)}
        \end{twocolentry}
        \vspace{0.1 cm}
        \begin{onecolentry}
            \begin{highlights}
                \item \textbf{Ascolto attivo e comunicazione empatica:} gestione quotidiana di richieste complesse, de-escalation di tensioni e risoluzione tempestiva con approccio orientato alla cura e alla chiarezza verso l'utente.
                \item \textbf{Gestione operativa e CRM:} utilizzo continuativo di piattaforme gestionali, tracciamento puntuale delle pratiche e precisione metodologica nell'organizzazione dei flussi di lavoro.
                \item \textbf{Chiarezza comunicativa:} attitudine a tradurre concetti tecnici complessi in spiegazioni semplici, accessibili e rassicuranti.
            \end{highlights}
        \end{onecolentry}

        \vspace{0.25 cm}

        \begin{twocolentry}{
            \textit{Catania, Italia}\vspace{0 cm}
            
            \textit{Precedente}
        }
            \textbf{Front-of-House \& Assistenza al Pubblico (Ristorazione)}\vspace{0 cm}
            
            \textit{Attività Ho.Re.Ca.}
        \end{twocolentry}
        \vspace{0.1 cm}
        \begin{onecolentry}
            \begin{highlights}
                \item Gestione del pubblico in contesti ad alta intensità, sviluppando resistenza allo stress, tempestività e cura dell'accoglienza.
                \item Comprensione pratica delle dinamiche alimentari reali, preferenze del consumatore e gestione di esigenze speciali (allergie, intolleranze, scelte etico-alimentari).
            \end{highlights}
        \end{onecolentry}

    \vspace{0.1 cm}

    % COMPETENZE & SUPPORTO ALLO STUDIO
    \section{Competenze per l'Affiancamento \& Aree di Interesse}
    \begin{itemize}[leftmargin=*,noitemsep,topsep=0pt]
        \item \textbf{Supporto Clinico ed Elaborazione Piani:} Rilevazione dati anamnestici e diari alimentari, calcolo fabbisogno energetico e macro (LARN, TDEE, BMR), stesura di tabelle di sostituzione alimentare, formulazione di ricette e grammature personalizzate.
        \item \textbf{Alimenti Vegetali \& Varietà Nutrizionale:} Conoscenza approfondita dei prodotti e degli alimenti di origine vegetale; capacità di comprendere le abitudini del paziente e consigliare opzioni pratiche e gradevoli da integrare nel piano alimentare, arricchendo la varietà della dieta con ascolto ed empatia.
        \item \textbf{Sensibilità Clinica \& DCA:} Approccio non giudicante al cibo, consapevolezza psiconutrizionale ed educazione alimentare orientata a sostenibilità e benessere duraturo.
        \item \textbf{Software \& Competenze Informatiche:} Suite Microsoft Office (Excel, Word), Google Workspace; rapido apprendimento nell'uso di software nutrizionali (Metadieta, WinFood, Nutrium).
    \end{itemize}

\end{document}
